\section{CVEX v3}
\label{sec:format}

\cvexthree has three defining properties. First, it is \emph{Docker-only}: a
record is a Docker Compose project, with no virtual machines, hypervisor, or
per-host provisioning. Second, every record is exercised twice, once with
legitimate activity and once with the attack, and both runs are captured as
network packets and system calls. Third, success is \emph{flag-based}: the
record plants a secret that only a successful exploitation can reach, and
the attack succeeds if and only if it retrieves it. The format descends from
the previous generations described in Section~\ref{sec:background}: it keeps
the capture sidecar and the explicit exploit verification of the first
generation, the server/client container pattern of the second, and the
\file{cvex.yml} manifest and process tracing of the VM-based framework. The
flag and the benign/attack pairing are new.

\subsection{Record Structure}

Every record has the same shape, so that tools and researchers can treat all
records uniformly. A record is a directory named after its CVE identifier:

\begin{lstlisting}
CVE-YYYY-NNNNN/
  cvex.yml            # manifest (Sec. 3.3)
  compose.yml         # derived Docker Compose project
  target/
    Dockerfile        # builds the vulnerable software
    vendor/           # pinned sources and binaries
    plant-flag.sh     # installs the flag (Sec. 3.4)
  benign/
    Dockerfile
    interact.sh       # legitimate use of the service
  attacker/
    Dockerfile
    exploit.sh        # the attack; must emit the flag
  recorder/           # shared capture image
  _source/            # provenance: original artifacts
  dataset/            # committed reference run (Sec. 3.7)
\end{lstlisting}

A record is composed of four container roles, each with a single
responsibility (Table~\ref{tab:roles}). A record may add a second target
container, typically a database, when the service needs one; everything else
stays fixed. The \file{\_source} directory preserves the artifacts from which
the record was derived (for example, the previous-generation record or the
\genie reproduction, together with the CVE metadata), so that the provenance
of every record is explicit.

\begin{table}[t]
\centering
\caption{Container roles in a \cvexthree record.}
\label{tab:roles}
\small
\begin{tabular}{@{}lp{0.68\columnwidth}@{}}
\toprule
\textbf{Role} & \textbf{Responsibility} \\
\midrule
\texttt{target} & runs the vulnerable software and holds the planted flag \\
\texttt{benign} & a client that uses the service as a legitimate user would \\
\texttt{attacker} & a client that exploits the service and retrieves the flag \\
\texttt{recorder} & captures packets and system calls during each phase \\
\bottomrule
\end{tabular}
\end{table}

The roles are arranged in one of three \emph{topologies}. In the
\texttt{target-attacker} topology, the benign and attacker clients are
separate containers that reach the target over a private network; in
\texttt{target+db-attacker}, the target is backed by a database container.
In the \texttt{single} topology, used for local vulnerabilities and for
libraries and command-line tools that have no network service, the benign
workload and the attack run inside the target container itself, as
\file{target/benign.sh} and \file{target/attack.sh}.

\subsection{Design Rationale}

Two choices shape the rest of the format. The first is that a single, uniform
notion of success is more valuable than a per-record test. A record whose
exploit ``ran'' tells a user very little: the exploit may have failed
silently, or the environment may not have been vulnerable at all. By
requiring every attack to produce a value that only a successful
exploitation can obtain, the format guarantees that the attack half of every
dataset is a genuine, successful exploitation, and makes verification a
simple string comparison that is identical across vulnerability classes.

The second choice is to separate the benign and attack phases rather than
mixing them. Running the two phases against identical, pristine targets
makes the two captures directly comparable, so that the difference between
them is the attack itself, and gives clean labels. A \emph{mixed} phase, in
which benign activity runs concurrently with the attack, can be added later
without changing the schema.

\subsection{The Manifest}

The \file{cvex.yml} manifest declares everything the runner needs: the
topology, the network layout, the flag and the success criterion, the
container images, the processes to trace, and the two phases.
Listing~\ref{lst:manifest} shows a manifest for a demonstration record that
serves a benign \texttt{nginx} instance.

\begin{lstlisting}[float=t,caption={A \file{cvex.yml} manifest (demonstration record).},label={lst:manifest}]
cvex_version: 3
cve: CVE-0000-00001
topology: target-attacker # | single | target+db-attacker
network:
  subnet: 10.13.0.0/24    # fixed addresses in the pcap
flag:
  value: "CVEX{0000000000000000000000000000demo}"
  plant: file:/opt/secret/flag
success: flag             # flag | verifier | crash
capture:
  flag_channel: network   # stdout | network | verifier
  tls_intercept: true
  tls_ports: [443]
services:
  target:
    image: cvex/cve-0000-00001-target@sha256:<digest>
    ip: 10.13.0.10
    trace: nginx          # process(es) to trace (regex)
    ready: "tcp:443"      # readiness condition
phases:
  benign:
    client: benign        # runs benign/interact.sh
    image: cvex/cve-0000-00001-benign@sha256:<digest>
    ip: 10.13.0.30
  attack:
    client: attacker      # runs attacker/exploit.sh
    image: cvex/cve-0000-00001-attacker@sha256:<digest>
    ip: 10.13.0.20
limits:
  wall_clock_seconds: 600
\end{lstlisting}

The \texttt{ready} condition tells the runner when the target is ready for a
phase to start: either a TCP port that accepts connections or a file that the
target creates when its initialization completes. The \texttt{trace} field is
a regular expression that selects the processes whose system calls are
recorded. Records converted from \genie additionally carry an
\texttt{origin} field and the upstream \texttt{source} and
\texttt{sw\_version} of the vulnerable software.

The format imposes a small set of rules, checked at review time: every image
is pinned by an \texttt{@sha256} digest; either \texttt{flag.value} and
\texttt{flag.plant} are present, or the record declares a different success
criterion; both phases are present; and addresses are static, so that
captures are reproducible.

\subsection{Success Criteria}

\paragraph{The flag.}
A flag is a short token of the form \texttt{CVEX\{}\textit{32 hex
digits}\texttt{\}}, planted inside the target at a location, or behind a
protection, that only the vulnerability can bypass. The secrecy of the flag
comes from \emph{access control}, not from its being unknown: a legitimate
user of the service cannot reach it, the benign client never sees it, and
only a successful exploitation returns it. Because the flag is gated by
access rather than by randomness, it can be a fixed value committed with the
record, which keeps the reference dataset byte-reproducible. A runner may
optionally generate a fresh flag for each run, to guard against an attack
that simply hard-codes the expected value; the gating is what proves the
exploitation either way.

The location of the flag follows the class of the vulnerability, using a
small set of standard placements (Table~\ref{tab:flags}). The
\texttt{plant} field of the manifest names the location and, where relevant,
the file mode, and \file{plant-flag.sh} installs the flag at container start.

\begin{table}[t]
\centering
\caption{Standard flag placements by vulnerability class.}
\label{tab:flags}
\small
\begin{tabular}{@{}p{0.36\columnwidth}p{0.56\columnwidth}@{}}
\toprule
\textbf{Class} & \textbf{Flag location} \\
\midrule
Command execution / injection & \file{/flag}, readable only by running a command on the target \\
Path traversal / file read & a file outside the served directory \\
File upload leading to execution & \file{/flag}, readable once uploaded content is executed \\
SQL injection & a table row that the application never exposes \\
Privilege escalation & \file{/root/flag}, owned by root, mode 0600 \\
Local code execution / deserialization & \file{/flag} or an environment value not exposed by the normal interface \\
Memory corruption (crash only) & none; a crash witness is used instead \\
\bottomrule
\end{tabular}
\end{table}

\paragraph{Capture channels.}
The flag returned by the attack is detected on the channel that the record
declares. With the \texttt{stdout} channel, the attack prints the flag, and
the runner extracts it from the attack's output; this is the universal
fallback, and the natural choice for local records. With the
\texttt{network} channel, the runner searches for the flag in the attack
packet capture, which also proves that the flag crossed the wire.

\paragraph{The verifier.}
Records converted from \genie use the \texttt{verifier} criterion. \genie
produces, for each reproduced CVE, a verifier program in the style of a CTF
checker: it sets up the vulnerable condition, runs the exploit, checks for
the effect of the vulnerability, and prints a fixed flag if and only if the
check passes. A \cvexthree record that declares \texttt{success: verifier}
runs this program inside the target during the attack phase, and the runner
accepts the reproduction if the verifier exits successfully and prints the
flag. The verifier criterion extends the flag idea to vulnerabilities whose
effect is not the disclosure of a secret, for example a parser that splits
one request into two.

\paragraph{The crash witness.}
Some vulnerabilities, such as memory-corruption bugs that are reproduced only
up to a crash, have no secret to read. For these, the record declares
\texttt{success: crash}, and the flag is replaced by a \emph{crash witness}:
the recorder confirms that a traced process terminated because of a signal,
or produced a core dump, during the attack phase and not during the benign
phase. This is the only deviation from the flag rule, and it is declared
explicitly.

In the current corpus, \NumFlag functional records use a planted flag
(\NumChanStdout on the \texttt{stdout} channel and \NumChanNetwork on the
\texttt{network} channel), \NumVerifier use a \genie verifier, and
\NumCrash uses a crash witness.

\subsection{Run Protocol}

A record is run by the reference runner in two isolated phases, each starting
from a \emph{pristine} target: the target containers are created fresh for
each phase, so that neither run can contaminate the other.

\begin{enumerate}
\item \textbf{Build.} The runner builds (or pulls) the images.
\item \textbf{Benign phase.} The runner starts the target (and database, if
any) and the recorder, plants the flag, waits for the \texttt{ready}
condition, starts the capture, runs the benign workload, stops the capture,
and tears the containers down.
\item \textbf{Attack phase.} The runner starts a fresh target and recorder,
plants the flag, waits for readiness, starts the capture, runs the attack,
collects the flag it emits, stops the capture, and tears the containers
down.
\item \textbf{Verify.} The runner compares the observed flag with the planted
one (or evaluates the verifier or the crash witness) and writes the dataset
and the verdict.
\end{enumerate}

The benign workload performs legitimate interactions with the same service
and, where possible, with the same endpoints that the attack uses: a normal
login, a normal search, a normal file request, a normal invocation of a tool.
The two captures are therefore directly comparable: same images, same
addresses, same target state at the start of the phase. Each phase is bounded
by the wall-clock limit declared in the manifest.

A second command reruns a record and compares the fresh dataset with the
committed one, checking that the flags match and comparing the traffic after
normalizing timestamps and ephemeral ports. Once a record has been verified in
this way, it is locked: a lock file records the hash of the verified record,
so that later modifications are detected.

\subsection{The Recorder}

The recorder is a single shared image. In \emph{sidecar} mode, it shares the
network namespace and the process namespace of the target container. It runs
\texttt{tcpdump} on the target's interfaces, including the loopback
interface, so that it also captures traffic between processes inside the
target, and it attaches \texttt{strace} to the processes selected by the
\texttt{trace} expression, following their children. In \emph{gateway} mode,
intended for multi-host records that use TLS, the recorder acts as the
default gateway of the record's network, captures traffic on the bridge, and
runs a transparent \texttt{mitmproxy} whose certificate authority is
installed in the client and target images at build time, so that encrypted
traffic can be decrypted reproducibly.

\subsection{Dataset Layout}
\label{sec:dataset}

The output of a run is a self-contained dataset directory:

\begin{lstlisting}
dataset/
  benign/
    benign.pcap          # legitimate traffic
    benign.syscalls.log  # system calls under normal use
  attack/
    attack.pcap          # attack traffic
    attack.syscalls.log  # system calls during the attack
  flag.observed          # the flag the attack retrieved
  verdict.json           # criterion, channel, planted vs.
                         # observed flag, success
  MANIFEST.sha256        # content hashes of all files
\end{lstlisting}

For local records the packet captures may be empty or contain only loopback
traffic, and the system-call traces carry most of the signal. The verdict
records the success criterion, the capture channel, the planted and the
observed flag, the outcome, and which captures are present.
Listing~\ref{lst:verdict} shows the verdict of the privilege-escalation
record discussed in Section~\ref{sec:examples}. The dataset is committed with
the record and addressed by the hashes in \file{MANIFEST.sha256}, so that the
reference captures survive even if the images can no longer be rebuilt.

\begin{lstlisting}[float=t,caption={The \file{verdict.json} of a functional record.},label={lst:verdict}]
{
  "cve": "CVE-2023-22809",
  "success": true,
  "criterion": "flag",
  "channel": "stdout",
  "flag_planted":  "CVEX{9b8bb059600aa99e8a28eca00bb4ddef}",
  "flag_observed": "CVEX{9b8bb059600aa99e8a28eca00bb4ddef}",
  "benign": { "pcap": true, "syscalls": true },
  "attack": { "pcap": true, "syscalls": true }
}
\end{lstlisting}

\subsection{Reproducibility}

The format pursues reproducibility at four levels. Builds are meant to be
offline and pinned: base images are referenced by digest, and the vulnerable
software is vendored into the record, so that the application layer can be
built without network access. The runtime is pinned: a static subnet and
static addresses, explicit readiness conditions, and a wall-clock limit. The
flag is fixed and committed, so that reference captures are stable while
still proving the exploitation. Finally, the reference dataset is committed
and content-addressed.

Not all of these properties hold yet for every record. Images are not yet
pinned by digest, and records converted from \genie still download the
vulnerable software at build time rather than vendoring it. Packet timing is
not reproducible in any case: payloads and the order of exchanges are, and
timing is normalized away when datasets are compared.

\subsection{Limitations}

Because containers share the host kernel, the format cannot reproduce
vulnerabilities in the kernel itself, or in operating systems other than
Linux; Section~\ref{sec:future} discusses an extension to virtual machines
that shares the same manifest and dataset format. TLS interception assumes
that the client and target trust the recorder's certificate authority, which
is acceptable for research reproductions and is noted in each dataset.
Finally, the realism of the benign phase is defined by the record's author:
the value of the benign baseline depends on how representative the benign
workload is, and reviewers should check that it exercises the same interface
as the attack.
